\section{Methods}
\label{sec:methods}

\subsection{Baseline feature engineering and ML pipeline}
\label{sec:methods-baseline}

Our baseline pipeline produces approximately 57 features per 1~Hz row,
grouped into six categories:

\begin{itemize}
  \item \textbf{Temporal:} \texttt{hour}, \texttt{dayofweek}, \texttt{month},
    \texttt{year}, \texttt{is\_weekend}, \texttt{is\_rush\_hour} (derived
    from the ISO timestamp).
  \item \textbf{Spatial:} latitude and longitude in degrees (converted from
    radians) and GNSS altitude.
  \item \textbf{Operational:} speed in km/h (from odometry), acceleration
    (first difference of speed), and rolling mean and standard deviation of
    speed over 60-second and 300-second windows. Rolling features are
    computed per-mission to prevent cross-mission bleeding.
  \item \textbf{Sensor:} electric power demand, traction force, brake
    pressure, and rolling mean and standard deviation of power over 60-second
    and 300-second windows.
  \item \textbf{Status:} binary flags for door open, grid available,
    halt brake active, and park brake active.
  \item \textbf{Categorical:} one-hot encoded \texttt{busRoute} (9 columns
    after drop-first) and top-20 \texttt{stopName} (20 columns + one
    \texttt{\_\_other\_\_} bucket).
\end{itemize}

We train three models on this feature set:
\begin{itemize}
  \item \textbf{Decision Tree}~\cite{pedregosa2011sklearn} (scikit-learn,
    Gini criterion, unlimited depth).
  \item \textbf{Random Forest}~\cite{breiman2001rf} (100 trees, parallel
    across CPUs, no max depth).
  \item \textbf{XGBoost}~\cite{chen2016xgboost} (500 trees, max depth 6,
    learning rate 0.05, histogram tree method).
\end{itemize}

All models are trained and evaluated independently under the pooled and
temporal frameworks. We report macro-F1 and balanced accuracy on the test
partition. Training uses all 47.6~M rows.

\subsection{Time series diagnostics}
\label{sec:methods-tsa}

Before introducing temporal features, we verified that the passenger count
series exhibits exploitable temporal structure at two granularities:
stop-event level within a mission (intra-trip autocorrelation) and hourly
aggregate level across all missions (inter-trip temporal structure).

\textbf{Stationarity tests.} We applied the Augmented Dickey-Fuller (ADF)
test~\cite{dickey1979unit} and the KPSS test~\cite{kwiatkowski1992kpss} in
parallel. ADF tests the null hypothesis of a unit root (non-stationarity);
rejection implies stationarity. KPSS tests the null of stationarity;
rejection implies non-stationarity. Using both tests in combination
provides robust evidence, as they have opposite nulls and complementary
statistical power.

For the hourly aggregate series, ADF rejects at $p < 0.001$ and KPSS rejects
at $p = 0.01$, a signature of \emph{trend-stationarity} (the series is
stationary around a trend, not around a constant mean). First differencing
produces a fully stationary series. For representative stop-level series
from individual missions, outcomes vary: three are stationary, one is
trend-stationary, one is inconclusive (neither test rejects).

\textbf{Autocorrelation diagnostics.} Figure~\ref{fig:acf} shows the
autocorrelation function (ACF) of the hourly aggregate series out to 48 lags
(two days). All 48 lags exceed the 95\% significance band, indicating
substantial temporal memory. Figure~\ref{fig:pacf} shows the corresponding
partial autocorrelation function (PACF): lag~1 is the dominant direct
effect ($\approx 0.58$), lag~2 shows a significant negative coefficient
($\approx -0.18$) consistent with mean reversion, and scattered significant
spikes extend out to lag~41. This suggests that an autoregressive model of
moderate order should capture the dominant temporal structure.

\begin{figure}[t]
  \centering
  \includegraphics[width=\columnwidth]{03_acf_hourly.png}
  \caption{Autocorrelation function (ACF) of the hourly aggregate passenger
    count series, out to 48 lags. All 48 lags exceed the 95\% significance
    band (shaded), indicating substantial temporal memory.}
  \label{fig:acf}
\end{figure}

\begin{figure}[t]
  \centering
  \includegraphics[width=\columnwidth]{05_pacf_hourly.png}
  \caption{Partial autocorrelation function (PACF) of the hourly aggregate
    series. Lag~1 dominates ($\approx 0.58$); lag~2 shows a significant
    negative effect consistent with mean reversion.}
  \label{fig:pacf}
\end{figure}

\subsection{ARIMA baseline (temporal-only floor)}
\label{sec:methods-arima}

As a sanity check that temporal structure alone is insufficient for
classification, we fit an $\text{ARIMA}(p, d, q)$ model to the hourly
aggregate passenger count series~\cite{box1976timeseries}. Using an AIC grid
search over $p \in \{0, \ldots, 10\}$, $q \in \{0, \ldots, 3\}$, and $d = 1$
(confirmed from stationarity analysis), the optimal order is
$\text{ARIMA}(8, 1, 3)$. We forecast the test period (2022+ hours),
threshold the continuous predictions using the training tercile boundaries
to produce class labels, and evaluate the resulting classification.

Figure~\ref{fig:arima} shows the forecast versus actual values.
Quantitatively, macro-F1 is 0.1795 and balanced accuracy is 0.335; the
confusion matrix reveals the model predicts ``medium'' for over 99\% of
test hours. This confirms that \emph{temporal signal alone is insufficient
for this classification task}. However, the ARIMA result also confirms that
\emph{temporal structure exists}, justifying the feature engineering approach
that follows.

\begin{figure*}[t]
  \centering
  \includegraphics[width=0.85\textwidth]{06_arima_baseline.png}
  \caption{$\text{ARIMA}(8, 1, 3)$ forecast of hourly aggregate ridership
    vs. actual values, 2022 test period (top). The model predicts a smooth
    near-mean sequence, collapsing to ``medium'' for over 99\% of test
    hours (confusion matrix, bottom) and yielding macro-F1 = 0.18.}
  \label{fig:arima}
\end{figure*}

\subsection{The forward-fill lag trap (v1, failed)}
\label{sec:methods-v1}

Motivated by the strong autocorrelation observed in Figures~\ref{fig:acf}
and~\ref{fig:pacf}, we added lag features to the baseline pipeline. Our
first implementation used the obvious approach: apply a \texttt{shift(N)}
operation to the forward-filled \texttt{itcs\_numberOfPassengers} column at
the 1~Hz granularity, producing features \texttt{pax\_lag\_60} (passenger
count 60 seconds ago) and \texttt{pax\_lag\_300} (passenger count 300
seconds ago). These lag sizes align with the existing 60-second and
300-second rolling window features.

The resulting model performance was surprisingly strong: Decision Tree
pooled macro-F1 jumped from 0.5713 (baseline) to 0.8743, a $+0.30$
improvement from two features alone on a balanced 3-class task. This is
suspicious by inspection, and we investigated.

\textbf{Diagnosis.} Recall that passenger count is constant between stop
events due to the forward-fill. Stop events occur roughly every 60--120
seconds. The \texttt{shift(60)} operation looks back exactly 60 rows, which
at 1~Hz is exactly 60 seconds. For a given row to have a different
\texttt{lag(60)} value than its current value, a stop event must have
occurred within the 60-second window $[t - 60, t)$. But most rows fall
entirely within a single inter-stop interval, so $\text{lag}(60) = p(t)$ for
the vast majority of rows. Empirically, we estimate this equality holds for
approximately 95\% of rows.

The model was not learning temporal structure; it was reading a near-copy
of the target variable. The Decision Tree effectively learned:
``if \texttt{pax\_lag\_60} $\le q_1$ then predict \emph{low}; else if
$\le q_2$ then predict \emph{medium}; else predict \emph{high}'', which
is the same rule as applying tercile boundaries directly to
\texttt{pax\_lag\_60}. Since $\text{pax\_lag\_60} \approx p(t)$ for 95\% of
rows, this rule achieves near-perfect accuracy on those rows.

\begin{figure}[t]
  \centering
  \includegraphics[width=\columnwidth]{fig_08_leakage_diagram.pdf}
  \caption{The forward-fill lag trap. Top: the raw passenger count (sparse
    event signal). Middle: the forward-filled dense series (staircase).
    Bottom: applying \texttt{shift(60)} at 1~Hz returns a nearly-identical
    copy of the middle series because the 60-second window rarely crosses
    a stop event. The overlap region (shaded) contains the leakage:
    \texttt{lag(60)} equals the current value for approximately 95\% of
    rows.}
  \label{fig:leakage}
\end{figure}

This is a target leakage pattern that survives every standard safeguard.
Mission-level splits are clean (no mission appears in both train and test),
cross-validation is honest, and no test-set data has entered the pipeline.
The leakage is \emph{endogenous to the feature construction}; it arises
from the interaction between imputation and feature computation, both of
which are individually benign.

\subsection{The fix: stop-level lag features (v2)}
\label{sec:methods-v2}

We redesigned the lag feature computation to operate at the natural event
granularity. Instead of shifting the forward-filled 1~Hz series, we shift
only the stop-event rows and then forward-fill the \emph{lag values} onto
the dense grid.

Formally, let $\mathcal{S} = \{t_1, t_2, \ldots, t_S\}$ be the set of
stop-event row indices in a mission, ordered chronologically, and $p(t)$
be the passenger count at row $t$. The naive (v1) approach computes
$\text{lag}_{60}^{\text{v1}}(t) = p(t - 60)$, which equals $p(t)$ whenever
$(t - 60, t) \cap \mathcal{S} = \emptyset$. Our corrected (v2) approach
computes, for each $k \in \{1, 3, 5\}$:
\begin{equation}
  \text{stoplag}_k(t) = p(t_{i - k})
\end{equation}
where $t_i$ is the most recent stop event on or before row $t$, i.e.,
$t_i = \max\{t_j \in \mathcal{S} \mid t_j \le t\}$. In words: at each stop
event $t_i$, we assign $\text{stoplag}_k$ the passenger count at
$t_{i-k}$, the passenger count $k$ stops ago, and then forward-fill
this value through all inter-stop rows until the next stop event updates
it.

We implement this in three steps:
\begin{enumerate}
  \item Mask the forward-filled series using \texttt{is\_stop\_event} to
    recover the sparse stop-only sequence.
  \item Apply \texttt{shift(k)} on the stop-only sequence.
  \item Reindex the resulting sparse lagged series to the full 1~Hz grid
    and forward-fill.
\end{enumerate}

By construction, $\text{stoplag}_k(t)$ is the passenger count at a strictly
earlier stop event. It cannot equal the current passenger count unless the
intervening stops happened to have the same count (which is rare in practice
and, when it does occur, reflects genuine short-term stability rather than
leakage). Figure~\ref{fig:stoplag} diagrams the computation.

\begin{figure}[t]
  \centering
  \includegraphics[width=\columnwidth]{fig_09_stop_lag_broadcast.pdf}
  \caption{Stop-level lag feature construction (v2). Top: the raw sparse
    event series. Middle: $\text{stoplag}_1$ assigned at each stop event
    as the previous stop's passenger count. Bottom: the lag values
    forward-filled onto the dense 1~Hz grid. Each row now has a lag feature
    that comes from a strictly earlier stop event and cannot equal the
    current count.}
  \label{fig:stoplag}
\end{figure}

We use $k \in \{1, 3, 5\}$, producing three stop-level lag features at
offsets of one, three, and five prior stop events. Lag~1 captures the
dominant short-range autocorrelation identified by the PACF analysis;
lags~3 and~5 provide longer-range context while remaining within a
typical trip.

\subsection{TimesFM embedding extraction}
\label{sec:methods-timesfm}

For our third feature set, we extract representations from a pretrained
time series foundation model. TimesFM~2.5-200M~\cite{das2024timesfm} is a
decoder-only transformer with 20 layers, hidden dimension 1280, and 16
attention heads, trained by Google Research on a large corpus of public
time series data.

\textbf{Architecture probe.} We verified the internal module structure of
the official PyTorch release of TimesFM~2.5-200M by registering forward
hooks and logging output shapes on a GPU compute node.
The computation flow is:
\begin{enumerate}
  \item Input time series passed through a patched tokenizer
    (a \texttt{ResidualBlock} with a linear projection).
  \item The resulting tokens pass through \texttt{stacked\_xf}, a
    \texttt{ModuleList} of 20 \texttt{Transformer} layers. Each layer
    contains RMSNorm + MultiHeadAttention (16 heads, 80 dims each) +
    RMSNorm + Feedforward (two linear layers with SiLU activation).
  \item The output of \texttt{stacked\_xf[19]} (the last transformer layer)
    passes to two projection heads:
    \texttt{output\_projection\_point} (point forecasts) and
    \texttt{output\_projection\_quantiles} (uncertainty bounds).
\end{enumerate}

At the output of the last transformer layer, before the projection heads
compress the information to forecast scalars, each token in the sequence
has a 1280-dimensional representation encoding the model's accumulated
understanding of the temporal context after 20 rounds of self-attention
and feedforward processing.

\textbf{Extraction pipeline.} For each of the 1,409 missions, we build
overlapping 512-second context windows of the forward-filled passenger
count with a 60-second stride. Each window is passed through the TimesFM
forecast pipeline, and a forward hook registered on \texttt{stacked\_xf[-1]}
captures the penultimate-layer output. We then mean-pool across the
sequence dimension to produce one 1280-dimensional vector per window, and
mean-pool across all windows within a mission to produce one
1280-dimensional vector per mission. Double mean-pooling is admittedly lossy,
but it produces a fixed-size per-mission representation suitable for use as
tabular features.

\textbf{PCA reduction.} The 1280-dimensional raw embeddings are reduced
via PCA to 32 dimensions. The PCA is fit \emph{exclusively on training-set}
mission embeddings (pre-2022) and then applied to all missions, preventing
test-set information from influencing the reduction. Remarkably, the 32
retained components capture \emph{100\% of the variance} in the training
embeddings, indicating that the foundation model's representations for
ZTBus bus ridership data live in a low-rank subspace of dimension at most
32. We return to this observation in Section~\ref{sec:discussion}.

\textbf{Feature attachment.} Each 1~Hz row in the training matrix inherits
the 32-dimensional embedding of its parent mission. We implement this via
a single \texttt{pd.merge} on a 1409-row lookup table rather than iterative
per-mission \texttt{.loc} assignment. The merge-based implementation
completes in 15 seconds on the full 47~M-row dataset, compared to an
estimated 55 minutes for the iterative approach, avoiding both
out-of-memory failures and unnecessary wall-clock time. This is an
implementation detail but an important one: naive implementations of the
broadcast pattern frequently cause OOM on large tabular datasets.

\subsection{Reproducibility}

All code is released as an open-source repository, including the full
preprocessing pipeline, feature engineering, model training, evaluation,
and figure generation. The pipeline includes 243 unit tests covering each
module, written before implementation (test-first development). All
stochastic operations are seeded for reproducibility. Cluster jobs use
checkpointed training with atomic JSON writes so that crashes and wall-time
kills can resume without re-computation.
